\begin{center}
\sc\large
Univerzita Komenského v Bratislave\\
Fakulta matematiky, fyziky a informatiky
\end{center}

\vspace{1cm}

\begin{center}
{\Large\bf Zadanie záverečnej práce}
\end{center}

\vspace{0.5cm}

\noindent
\begin{tabular}{@{}p{4.5cm}p{10cm}}
\textbf{Meno a priezvisko študenta:} & Vladyslav Havriuk \\[0.3cm]
\textbf{Študijný program:} & informatika (Jednoodborové štúdium, bakalársky I. st., denná forma) \\[0.3cm]
\textbf{Študijný odbor:} & informatika \\[0.3cm]
\textbf{Typ záverečnej práce:} & bakalárska \\[0.3cm]
\textbf{Jazyk záverečnej práce:} & anglický \\[0.3cm]
\textbf{Sekundárny jazyk:} & slovenský \\[0.3cm]
\end{tabular}

\vspace{0.5cm}

\noindent
\textbf{Názov:} Protocol fuzzing

\noindent
Fuzzing protokolov

\vspace{0.5cm}

\noindent
\textbf{Anotácia:}

\noindent
Protocol fuzzing is a technique used to test protocol implementations. It has been successfully employed to identify implementation errors and security vulnerabilities in various protocols. The goal of this thesis is to explore methods of protocol fuzzing and apply them to real-world implementations of the WebTransport protocol. The results will be analyzed and discussed.

\vspace{1cm}

\noindent
\begin{tabular}{@{}p{4.5cm}p{10cm}}
\textbf{Vedúci:} & doc. RNDr. Martin Stanek, PhD. \\[0.3cm]
\textbf{Katedra:} & FMFI.KI -- Katedra informatiky \\[0.3cm]
\textbf{Vedúci katedry:} & prof. RNDr. Martin Škoviera, PhD. \\[0.3cm]
\textbf{Dátum zadania:} & 23. 10. 2025 \\[0.3cm]
\textbf{Dátum schválenia:} & 27. 10. 2025 \\[0.3cm]
\end{tabular}

\vspace{1.5cm}

\noindent
\begin{tabular}{@{}p{7cm}p{7cm}}
\centering\rule{5cm}{0.4pt} & \centering\rule{5cm}{0.4pt}\tabularnewline
\centering študent & \centering garant študijného programu
\end{tabular}
